\chapter{The Subaltern Perspective}

\section{What is the Subaltern Perspective?}

The subaltern perspective means to \textbf{read history from below} --- from the perspective of the led, not the leaders. Nationalist historiography projects Gandhi, Nehru, Tilak etc.\ as leaders of peasant and tribal movements (though they were neither peasant nor tribal), making history \textbf{hero-centric}. The ignored movements and people are the subaltern.

\begin{ExampleBox}
``I fought with them'' vs ``they fought with him'' --- the difference is only the \textbf{camera angle}. Nationalist history keeps the camera on the hero; subaltern history turns it to the movement and the masses (e.g.\ what did the Champaran people think of Gandhi?).
\end{ExampleBox}

\begin{KeyConcept}
\textbf{Subaltern} literally means \textbf{inferior ranks} (not class --- ``class'' is Marxist terminology). It includes \textbf{peasants, tribals, Dalits}, and Dalit/tribal/peasant \textbf{women} (intersectional). The term was coined by \textbf{Antonio Gramsci} (a neo-Marxist): cultural hegemony vs \textbf{counter-hegemony} (the subaltern represents counter-hegemony).
\end{KeyConcept}

\begin{DataFact}
Subaltern is the \textbf{only perspective coming from a colonized country} (others come from colonizers); it is also called \textbf{post-colonial historiography}.
\end{DataFact}

\begin{QuickRevision}
\begin{itemize}
\item Subaltern = read history from below (camera angle on the movement, not the leader).
\item Inferior ranks: peasants, tribals, Dalits, their women (Gramsci coined the term; counter-hegemony).
\item Only perspective from a colonized country; = post-colonial historiography.
\end{itemize}
\end{QuickRevision}

\section{David Hardiman \& The Coming of Devi}

\textbf{David Hardiman} (a socialist historian) wrote \emph{The Coming of Devi} (Devi = Goddess), about an indigenous, autonomous tribal movement in the coastal belt of \textbf{Gujarat (1922--23)}, the \textbf{Devi movement}.

Parsi businessmen (liquor shops) + landlords + moneylenders formed a nexus that addicted tribals to alcohol, pushing them into debt and marginalisation. When \textbf{smallpox} struck in 1922, a rumour spread; \textbf{shamans} (believed possessed by the Goddess) organised rituals, and the Goddess ``said'': stop drinking alcohol. The Adivasis became teetotalers.

\begin{CriticalPoint}
This was an \textbf{autonomous} social movement with elements of \textbf{social reform}, not spearheaded by any elite nationalist bourgeois leader. Gandhi later incorporated Adivasis into the mass movement only because these local movements showed tribals were assertive.
\end{CriticalPoint}

\begin{QuickRevision}
\begin{itemize}
\item Hardiman, The Coming of Devi: autonomous Adivasi movement in coastal Gujarat (1922--23).
\item Nexus (Parsi liquor shops + landlords + moneylenders) leading to addiction, then smallpox, then Devi says stop alcohol, then teetotalism.
\end{itemize}
\end{QuickRevision}

\section{Ranajit Guha: Critique of Nationalist Historiography}

\textbf{Ranajit Guha} is the founder of Subaltern Studies. His critique of mainstream nationalist historiography:
\begin{itemize}
\item Nationalist historiography does not represent history --- it \textbf{produces heroes}.
\item The subalterns are not objects of history, but the \textbf{true makers of history}.
\item It focuses on peasants and tribals in the colonial movement, overlooked by elitist historiography.
\item Elitist historiography gives the impression that nationalist consciousness was the achievement of only the elites.
\item It constructs a \textbf{binary of elite and subaltern} (like Marx's bourgeoisie--proletariat).
\item Subaltern politics is an \textbf{autonomous domain}, independent of elite politics.
\end{itemize}

\begin{KeyConcept}
Guha's conclusion: the history of Indian nationalism is a \textbf{``spiritual biography of Indian elites''} --- we read biographies of leaders, not movements.
\end{KeyConcept}

\begin{ExampleBox}
The NCERT \emph{Themes in Indian History (Part 3)} reflects the subaltern turn: the \textbf{Pahadiyas} (hunters, shifting cultivators, charcoal producers) --- the backbone of the Santhal rebellion --- get attention; \textbf{Siddhu Manji} (Santhal leader); moneylenders projected as ``Dekus'' (dacoits). \textbf{Shahid Amin} shows the ``Mahatma'' title was first used by Gorakhpur people, later attributed to Tagore.
\end{ExampleBox}

\begin{QuickRevision}
\begin{itemize}
\item Ranajit Guha = founder of Subaltern Studies; nationalist history produces heroes, not history.
\item History of Indian nationalism = spiritual biography of elites.
\item Pahadiyas (backbone of Santhal rebellion), Siddhu Manji, Shahid Amin (Mahatma title).
\end{itemize}
\end{QuickRevision}

\section{B. R. Ambedkar \& the Subaltern}

Ambedkar is not a Marxist (though influenced by Marxism --- his book \emph{Marx and Buddha}; Buddha was against caste). The Communist Party fought for workers' rights, never for the \textbf{human rights of Dalits}.

\begin{KeyConcept}
Ambedkar: \textbf{Hinduism and democracy cannot coexist} --- democracy is equality, Hinduism is \textbf{graded inequality}. Caste is not just a division of labour but a \textbf{division of labourers}.
\end{KeyConcept}

\begin{itemize}
\item The Indian nationalist movement was a \textbf{Hindu-Brahminical upper-caste men's movement}, not an Indian one. Congress was a \textbf{Brahminical party}; its demands (Indianization of civil services, modern education, representation in assemblies) served upper-caste Brahmin men.
\item The word \textbf{Dalit} was coined by \textbf{Jyotiba Phule} (in \emph{Gulamgiri}). Ambedkar was influenced by Phule (British Raj better than Peshwa Raj --- meritocracy vs casteism).
\item Vedas, Smritis and Shastras justify Dalit oppression (like Marx's religion opiating the masses).
\item Ambedkar's programmes: \textbf{Mahad Satyagraha} (temple/water entry), the \textbf{Republican Party of India} (1936, aiming at a \textbf{caste-less}, not just classless, India), journals (Mook Nayak, Equality Janta, Bahishkrit Bharat), and the cause of peasant castes (\textbf{Mahar} and \textbf{Kunbi}).
\end{itemize}

\begin{CriticalPoint}
Ambedkar's conclusion: \textbf{India has not gained independence} --- upper-caste Brahmin men gained independence from the British; Dalits are still waiting for independence from the upper-caste Brahmin men.
\end{CriticalPoint}

\begin{QuickRevision}
\begin{itemize}
\item Ambedkar is not a Marxist; CPI ignored Dalit human rights.
\item Hinduism + democracy cannot coexist; caste = division of labourers.
\item Congress = Brahminical party; Dalit word coined by Phule; Republican Party (1936) aimed at caste-less India.
\end{itemize}
\end{QuickRevision}

\section{Great Tradition vs Little Tradition}

A binary again: \textbf{great tradition} (national, elite, urban scale) vs \textbf{little tradition} (local, folk, common masses).

\begin{ExampleBox}
The Ramayana has \textbf{400+ versions}. The television Ramayana is the elitist/great-tradition version; in the North East Ramlila is dance-based, Ram is shown as an avatar of Vishnu; some versions show Ram promoting untouchability (the Shambuka episode); one Indian village worships Ravan and does not celebrate Diwali. \textbf{Karwa Chauth} was a Rajasthan/Gujarat folk tradition that became a great tradition through Bollywood.
\end{ExampleBox}

\begin{QuickRevision}
\begin{itemize}
\item Great tradition (national/elite) vs little tradition (local/folk).
\item Ramayana has 400+ versions; Karwa Chauth went from folk to great tradition.
\end{itemize}
\end{QuickRevision}

\section{The Critique of the Subaltern Perspective}

\begin{itemize}
\item \textbf{Sumit Sarkar} (a subaltern historian) feuded with Ranajit Guha: the voices now represented in subaltern literature are not really subaltern but elites (e.g.\ upper-caste women).
\item \textbf{Gayatri Spivak} --- ``Can the Subaltern Speak?'' --- when upper-caste Brahmins speak on behalf of the subaltern, reality is distorted; the social location of the researcher matters.
\item \textbf{Kumkum Roy} (JNU) --- where are the \textbf{Vedic dasis} (female slaves)?
\item Who is subaltern? Dalits include elites (e.g.\ Mayawati --- a Dalit elite).
\end{itemize}

\begin{QuickRevision}
\begin{itemize}
\item Critique: Sumit Sarkar (elites represented as subaltern), Gayatri Spivak (Can the Subaltern Speak?), Kumkum Roy (Vedic dasis).
\item Problem of representation and of who is truly subaltern (Dalit elites).
\end{itemize}
\end{QuickRevision}
